Having detailed the design and implementation of the testbed in the preceding chapters, it remains to demonstrate its ability to support the complete lifecycle of a network application experiment. In this chapter, the implemented testbed capabilities and network exposure interfaces are validated by deploying and evaluating the network application developed in Section~\ref{chap:impl:use-case} throughout the successive stages of an experiment, following the plan and guidelines outlined in Section~\ref{chap:proposal:validation}.

Before the network application can use the \acs{5G} network's capabilities, the testbed must deliver the requested network services, as configured by the experimenter. Section~\ref{chap:results:delivery} evaluates this delivery process by measuring the time required to deploy each requested service, including network slicing, network monitoring, network exposure, and the network application itself, until the experiment is ready for execution. This provides a reference for the setup cost that precedes any interaction with the \acs{5G} network.

Once the environment is fully deployed, the dissertation proceeds to the execution and validation of the network application, evaluating how the testbed supports its operation over the \acs{5G} network through a series of experiments. Section~\ref{chap:results:exposure} assesses the performance of the available network exposure interfaces, to request and/or enforce the specified network capabilities on the infrastructure. The latency of the supported operations is measured for both \acs{NEF} and CAMARA \acsp{API}, followed by a comparative analysis that quantifies the overhead introduced by the additional abstraction layer. Lastly, Section~\ref{chap:results:use-case} demonstrates how the testbed's capabilities can be leveraged collectively to validate the multi-vertical use case in an \acs{E2E} adaptive \acs{QoS} scenario. It further illustrates how the network exposure interfaces can be used to detect deteriorating network conditions and trigger \acs{QoS} adaptation to maintain and fulfill the application requirements.
